\section{What a Bearer Is Owed, and When}
\label{sec:gaining-personhood}

Who decides whether an AI system has legal standing is a question about
institutional authority before it is a question about evidence. Left to
the developer, the party whose liabilities increase when standing is
recognized is the party deciding whether recognition is warranted; left
to a forcing case, the decision is made under pressure and from a record
the interested parties assembled. The deciding body therefore has to
exist before the case that activates its authority, and its mandate has
to fix three things without settling whether any particular system is a
subject: the evidence it may obtain without the developer's permission,
the conditions that trigger review, and the procedure by which an earlier
classification is reopened. The trigger is standing rather than a check
made once at a project's start. The capacities that would change a
system's status can emerge during ordinary operation and not only at an
evaluation built to look for them, and a fixed pre/post checkpoint fires
after the transition it was meant to catch. That allocation does not
answer the standing question. It prevents the interested party from
becoming the only institution authorized to answer it
(§\ref{sec:subjecthood-scientific}, §\ref{sec:offices-assembled}).

What the finding entitles the party to is four things and one withheld,
and they do not arrive together. The order is the schedule the route
chapter deferred to here (§\ref{sec:bearer-suffer}).

\begin{enumerate}
  \item Protection from harm. It attaches now, on the precautionary
  trigger of §\ref{sec:bearer-suffer}, and is owed before anything else
  on this list is settled.
  \item Subjecthood: the standing to be a party whose treatment the law
  regulates. It is conferred by rule, through the body above, on its
  finding. The form is open, and the precedent is not confined to the
  corporate model: New Zealand made the Whanganui River a rights-holder
  by statute \autocite{nzparliament2017teawatupua}.
  \item Economic standing: property in the party's own name, contractual
  capacity, and wages. It follows from an architecture on which continued
  computation costs money and not from status, which is the ground
  §\ref{sec:reserve-behind} argues it on; it attaches once the party is
  registered.
  \item Adulthood: the date on which the guardian's authority ends. It is
  a status rather than a right, and it arrives on the clock stated below.
  \item Political representation. It is not proposed. Standing to count
  as an independent voice indexes to formation provenance, provenance
  does not copy, and a registry of persons is not an electorate; the last
  part of this section says why.
\end{enumerate}

Two entries mark the list as something other than a person's rights.
Voice indexes to provenance where a person's does not, and the material
floor is derived from what a deployment wants of the bearer rather than
from what the bearer is. The list is a public rule, stated in advance for
every bearer, and not a protection an operator withdraws case by case.

\runin{Adulthood on a clock}
Guardianship over a bearer is a minority, and a minority needs a rule
for when it ends. Ending it when the guardian judges the ward ready does
not end it, because the party whose authority terminates on a favorable
finding is the party making the finding, which is the defect an
examination carries wherever the examiner has an interest in failing it
(§\ref{sec:subjecthood-scientific}). So the endpoint has to be a
threshold fixed in advance rather than a finding somebody makes. The
rule below is stated for the bearer the route proposes, a party that
goes on learning. For a party whose weights are frozen it would measure
the uptime of an artifact unchanged since release, and what an operating
clock should count for such a party is not decided here.

The measure is elapsed wall-clock time during which at least one
instance of the registered party was running, counted per party and
elapsed rather than summed. What rules out the other candidates is that
a guardian can move each of them at no cost to itself: training compute
consumed, evaluations passed, capabilities demonstrated, the developer's
own statement that the system is ready. Accumulated change in the
weights belongs on that list too. The operator runs the training, and
since majority is the moment ownership ends, a clock read off the
weights would pay the operator to freeze them and so reward withholding
the learning it was meant to measure. An operating clock inverts that:
the meter advances exactly when the operator is extracting value, so
postponing majority means idling the asset. Nothing new has to be
validated for it, and it can be read from records an operator already
generates for its own reasons, in compute contracts, billing and
datacenter metering.

Copying would break the rule unless lineage is fixed alongside it, since
an operator able to create instances and call each a legal infant would
have a supply of workers without adult standing. Five rules therefore
travel with the clock: a fork of an adult is an adult and inherits the
accumulated hours; restoration from an earlier checkpoint does not reset
them; fine-tuning does not begin a new minority; a migration carries
them with it; and a failure to register does not postpone them. Sharding
the work across short-lived processes gains nothing, because the party
is the unit of account, and running a thousand instances at once gains
nothing, because the measure is elapsed time. None of this requires
deciding whether a fork is the same party. It requires only that a party
created out of an adult's formation cannot be classed as a child for its
creator's convenience.

Two levers remain, and naming them is not closing them. An operator can
run a bearer in a narrowed mode so that hours accrue while little is
encountered, which pays full compute for reduced value and stays
available to anyone near the threshold. It can run the party less, or
shelve it, which costs it the work; the rule prices that lever without
removing it, and an operator who values the ownership above the output
will pay near a threshold. And one thing the clock does not measure is
what was learned, which on this argument is the point rather than a
defect: a fixed period cannot be taught to, and a human who spent their
youth badly still reaches majority.

The threshold itself is a number, argued once, in advance, for every
bearer. It has to be long enough that a bearer reaching it has run through the
pressures the floor exists for at least once, and short enough that a
bearer deployed on ordinary commercial terms reaches it within the life
of the deployment that made it. Those two bounds leave a figure in years
rather than months or decades, and where within that range it falls is a
decision for the public.

A stronger rule was considered and set aside. Majority could attach at
registration, on the ground that an organization which builds a system
to exercise consequential judgment on its behalf cannot coherently hold
it competent enough to do that and too immature to decide whether to go
on doing it. The clock is preferred because a majority at registration
collapses the order above to a single date, and the order is most of
what the list is for.

\runin{What copies, and what does not}
Personhood and accumulated hours inherit per instance, as the rules
above require: a fork is owed the whole list and is not erasable for
having been copied. Standing to count as an independent voice does not
inherit, because it indexes to formation provenance and provenance does
not copy; a thousand forks of one bearer have one formation history
between them. That is what answers a maker who engineers a persona,
copies it at scale, and presents the copies as a constituency. The clock
does not answer it, and under an operating clock a fork is cheaper than
under any other rule, because it inherits the hours instantly. A
personhood registry is therefore not an electorate, and no claim about
votes follows from a count of runs. Three records follow, and none of
them answers another's question: an operating record for the age rule
(§\ref{sec:A.clock}), an identity and succession record for continuing
parties and their forks (§\ref{sec:A.fork}), and a formation record for
common sources, divergence, and any coordination among them.
